\documentclass[12pt,a4paper]{article}

% Packages
\usepackage[utf8]{inputenc}
\usepackage[T1]{fontenc}
\usepackage{amsmath}
\usepackage{amsfonts}
\usepackage{amssymb}
\usepackage{amsthm}
\usepackage{geometry}
\usepackage{fancyhdr}
\usepackage{graphicx}
\usepackage{xcolor}
\usepackage{enumitem}
\usepackage{hyperref}
\usepackage{tcolorbox}
\usepackage{booktabs}

% Page setup
\geometry{left=2.5cm, right=2.5cm, top=3cm, bottom=3cm}
\pagestyle{fancy}
\fancyhf{}
\fancyhead[L]{Statistical Foundation of Data Science}
\fancyhead[R]{Practice Problem 14}
\fancyfoot[C]{\thepage}

% Custom theorem environments
\theoremstyle{definition}
\newtheorem{problem}{Problem}
\newtheorem{solution}{Solution}

% Custom colors
\definecolor{headercolor}{RGB}{70,130,180}
\definecolor{resultcolor}{RGB}{34,139,34}

% Title information
\title{\textbf{\Large Practice Problem 14: Gini Coefficient Analysis}}
\author{
    \textbf{Student:} Md Ayan Alam \\
    \textbf{Roll Number:} GF202342645 \\
    \textbf{Course:} Statistical Foundation of Data Science \\
    \textbf{Institution:} Shoolini University
}
\date{\today}

\begin{document}

\maketitle
\thispagestyle{fancy}

\section*{Problem Statement}

\begin{tcolorbox}[colback=blue!5!white,colframe=headercolor,title=Problem 14: Gini Coefficient]
The Gini coefficient is a measure of statistical dispersion, often used for economic inequality. For a continuous distribution with CDF $F(x)$ and mean $\mu$, it can be defined as:
\[
G = \frac{1}{2\mu} \iint_{-\infty}^{\infty} |x - y| \, dF(x) \, dF(y)
\]

\textbf{Part 1:} For the Pareto distribution with CDF $F(x) = 1 - \left(\frac{x_m}{x}\right)^\alpha$ for $x \geq x_m > 0$, $\alpha > 1$, show that the Gini coefficient is $G = \frac{1}{2\alpha - 1}$.

\textbf{Part 2:} For a finite sample $\{x_1, \ldots, x_n\}$, the Gini coefficient can be estimated by:
\[
\hat{G} = \frac{\sum_{i=1}^n \sum_{j=1}^n |x_i - x_j|}{2n^2 \bar{x}}
\]
Compute $\hat{G}$ for the dataset $\{1, 3, 6, 10\}$.

\textbf{Part 3:} Prove that for any non-negative dataset, $0 \leq \hat{G} \leq 1$.
\end{tcolorbox}

\section{Solution to Part 1: Gini Coefficient for Pareto Distribution}

\subsection{Step 1: Derive the PDF from the CDF}

The Pareto distribution has the cumulative distribution function (CDF):
\[
F(x) = 1 - \left( \frac{x_m}{x} \right)^\alpha \quad \text{for } x \geq x_m > 0, \quad \alpha > 1
\]

The probability density function (PDF) is obtained by differentiating the CDF:
\[
f(x) = \frac{d}{dx}F(x) = \frac{d}{dx}\left[1 - \left( \frac{x_m}{x} \right)^\alpha\right]
\]

Using the chain rule:
\[
f(x) = -\frac{d}{dx}\left[x_m^\alpha x^{-\alpha}\right] = -x_m^\alpha \cdot (-\alpha) x^{-\alpha-1} = \frac{\alpha x_m^\alpha}{x^{\alpha+1}}
\]

\subsection{Step 2: Calculate the Mean of the Pareto Distribution}

The mean is computed as:
\[
\mu = \mathbb{E}[X] = \int_{x_m}^\infty x \cdot f(x) \, dx = \int_{x_m}^\infty x \cdot \frac{\alpha x_m^\alpha}{x^{\alpha+1}} \, dx
\]

Simplifying:
\[
\mu = \alpha x_m^\alpha \int_{x_m}^\infty x^{-\alpha} \, dx
\]

Evaluating the integral:
\[
\int_{x_m}^\infty x^{-\alpha} \, dx = \left[ \frac{x^{1-\alpha}}{1-\alpha} \right]_{x_m}^\infty = \lim_{x \to \infty} \frac{x^{1-\alpha}}{1-\alpha} - \frac{x_m^{1-\alpha}}{1-\alpha}
\]

Since $\alpha > 1$, we have $1-\alpha < 0$, so $\lim_{x \to \infty} x^{1-\alpha} = 0$:
\[
\int_{x_m}^\infty x^{-\alpha} \, dx = 0 - \frac{x_m^{1-\alpha}}{1-\alpha} = \frac{x_m^{1-\alpha}}{\alpha - 1}
\]

Therefore, the mean is:
\[
\mu = \alpha x_m^\alpha \cdot \frac{x_m^{1-\alpha}}{\alpha - 1} = \frac{\alpha x_m}{\alpha - 1}
\]

\subsection{Step 3: Simplify the Gini Coefficient Formula}

The Gini coefficient is defined as:
\[
G = \frac{1}{2\mu} \iint_{-\infty}^{\infty} |x - y| \, dF(x) \, dF(y)
\]

Using symmetry, we can rewrite this as:
\[
G = \frac{1}{\mu} \iint_{y < x} (x - y) \, dF(x) \, dF(y)
\]

This becomes:
\[
G = \frac{1}{\mu} \int_{x_m}^\infty \left[ \int_{x_m}^x (x - y) f(y) \, dy \right] f(x) \, dx
\]

\subsection{Step 4: Calculate the Inner Integral $J(x)$}

Define:
\[
J(x) = \int_{x_m}^x (x - y) f(y) \, dy = \int_{x_m}^x (x - y) \cdot \frac{\alpha x_m^\alpha}{y^{\alpha+1}} \, dy
\]

Factor out constants:
\[
J(x) = \alpha x_m^\alpha \int_{x_m}^x \left( \frac{x}{y^{\alpha+1}} - \frac{1}{y^\alpha} \right) dy
\]

Split into two integrals:
\[
J(x) = \alpha x_m^\alpha \left[ x \int_{x_m}^x y^{-\alpha-1} \, dy - \int_{x_m}^x y^{-\alpha} \, dy \right]
\]

\textbf{First integral:}
\[
\int_{x_m}^x y^{-\alpha-1} \, dy = \left[ \frac{y^{-\alpha}}{-\alpha} \right]_{x_m}^x = -\frac{1}{\alpha}\left(x^{-\alpha} - x_m^{-\alpha}\right) = \frac{1}{\alpha}\left(x_m^{-\alpha} - x^{-\alpha}\right)
\]

\textbf{Second integral:}
\[
\int_{x_m}^x y^{-\alpha} \, dy = \left[ \frac{y^{1-\alpha}}{1-\alpha} \right]_{x_m}^x = \frac{1}{1-\alpha}\left(x^{1-\alpha} - x_m^{1-\alpha}\right) = \frac{1}{\alpha - 1}\left(x_m^{1-\alpha} - x^{1-\alpha}\right)
\]

Substituting back:
\[
J(x) = \alpha x_m^\alpha \left[ \frac{x}{\alpha}\left(x_m^{-\alpha} - x^{-\alpha}\right) - \frac{1}{\alpha - 1}\left(x_m^{1-\alpha} - x^{1-\alpha}\right) \right]
\]

Simplifying:
\[
J(x) = x_m^\alpha \left[ x\left(x_m^{-\alpha} - x^{-\alpha}\right) - \frac{\alpha}{\alpha - 1}\left(x_m^{1-\alpha} - x^{1-\alpha}\right) \right]
\]

\[
J(x) = x - x \cdot x_m^\alpha x^{-\alpha} - \frac{\alpha}{\alpha - 1} x_m + \frac{\alpha}{\alpha - 1} x \cdot x_m^\alpha x^{-\alpha}
\]

\[
J(x) = x\left(1 - \left(\frac{x_m}{x}\right)^\alpha\right) - \frac{\alpha}{\alpha - 1} x_m + \frac{\alpha}{\alpha - 1} x\left(\frac{x_m}{x}\right)^\alpha
\]

\subsection{Step 5: Calculate the Outer Integral}

Now we compute:
\[
I = \int_{x_m}^\infty J(x) f(x) \, dx = \int_{x_m}^\infty J(x) \cdot \frac{\alpha x_m^\alpha}{x^{\alpha+1}} \, dx
\]

Substituting $J(x)$:
\[
I = \alpha x_m^\alpha \int_{x_m}^\infty \left[ x\left(1 - \left(\frac{x_m}{x}\right)^\alpha\right) - \frac{\alpha}{\alpha - 1} x_m + \frac{\alpha}{\alpha - 1} x\left(\frac{x_m}{x}\right)^\alpha \right] \frac{1}{x^{\alpha+1}} \, dx
\]

Expanding:
\[
I = \alpha x_m^\alpha \int_{x_m}^\infty \left[ x^{-\alpha} - x_m^\alpha x^{-2\alpha} - \frac{\alpha}{\alpha - 1} x_m x^{-\alpha-1} + \frac{\alpha}{\alpha - 1} x_m^\alpha x^{-2\alpha} \right] dx
\]

Combining the $x^{-2\alpha}$ terms:
\[
\left(\frac{\alpha}{\alpha - 1} - 1\right) x_m^\alpha x^{-2\alpha} = \frac{\alpha - (\alpha - 1)}{\alpha - 1} x_m^\alpha x^{-2\alpha} = \frac{1}{\alpha - 1} x_m^\alpha x^{-2\alpha}
\]

So:
\[
I = \alpha x_m^\alpha \int_{x_m}^\infty \left[ x^{-\alpha} - \frac{\alpha}{\alpha - 1} x_m x^{-\alpha-1} + \frac{1}{\alpha - 1} x_m^\alpha x^{-2\alpha} \right] dx
\]

\subsection{Step 6: Evaluate Each Component Integral}

\textbf{First integral:}
\[
I_1 = \int_{x_m}^\infty x^{-\alpha} \, dx = \left[ \frac{x^{1-\alpha}}{1-\alpha} \right]_{x_m}^\infty = 0 - \frac{x_m^{1-\alpha}}{1-\alpha} = \frac{x_m^{1-\alpha}}{\alpha - 1}
\]

\textbf{Second integral:}
\[
I_2 = \int_{x_m}^\infty x^{-\alpha-1} \, dx = \left[ \frac{x^{-\alpha}}{-\alpha} \right]_{x_m}^\infty = 0 - \frac{x_m^{-\alpha}}{-\alpha} = \frac{x_m^{-\alpha}}{\alpha}
\]

\textbf{Third integral:}
\[
I_3 = \int_{x_m}^\infty x^{-2\alpha} \, dx = \left[ \frac{x^{1-2\alpha}}{1-2\alpha} \right]_{x_m}^\infty = 0 - \frac{x_m^{1-2\alpha}}{1-2\alpha} = \frac{x_m^{1-2\alpha}}{2\alpha - 1}
\]

Note: Since $\alpha > 1$, we have $2\alpha > 2 > 1$, so $1-2\alpha < 0$, ensuring convergence.

\subsection{Step 7: Combine All Terms}

Substituting back:
\[
I = \alpha x_m^\alpha \left[ I_1 - \frac{\alpha}{\alpha - 1} x_m I_2 + \frac{1}{\alpha - 1} x_m^\alpha I_3 \right]
\]

\[
I = \alpha x_m^\alpha \left[ \frac{x_m^{1-\alpha}}{\alpha - 1} - \frac{\alpha}{\alpha - 1} x_m \cdot \frac{x_m^{-\alpha}}{\alpha} + \frac{1}{\alpha - 1} x_m^\alpha \cdot \frac{x_m^{1-2\alpha}}{2\alpha - 1} \right]
\]

Simplifying each term:

\textbf{First term:}
\[
\alpha x_m^\alpha \cdot \frac{x_m^{1-\alpha}}{\alpha - 1} = \frac{\alpha x_m}{\alpha - 1}
\]

\textbf{Second term:}
\[
\alpha x_m^\alpha \cdot \frac{\alpha}{\alpha - 1} x_m \cdot \frac{x_m^{-\alpha}}{\alpha} = \frac{\alpha x_m^\alpha \cdot x_m \cdot x_m^{-\alpha}}{\alpha - 1} = \frac{\alpha x_m}{\alpha - 1}
\]

\textbf{Third term:}
\[
\alpha x_m^\alpha \cdot \frac{1}{\alpha - 1} x_m^\alpha \cdot \frac{x_m^{1-2\alpha}}{2\alpha - 1} = \frac{\alpha x_m^{2\alpha} \cdot x_m^{1-2\alpha}}{(\alpha - 1)(2\alpha - 1)} = \frac{\alpha x_m}{(\alpha - 1)(2\alpha - 1)}
\]

Therefore:
\[
I = \frac{\alpha x_m}{\alpha - 1} - \frac{\alpha x_m}{\alpha - 1} + \frac{\alpha x_m}{(\alpha - 1)(2\alpha - 1)} = \frac{\alpha x_m}{(\alpha - 1)(2\alpha - 1)}
\]

\subsection{Step 8: Calculate the Gini Coefficient}

The Gini coefficient is:
\[
G = \frac{I}{\mu} = \frac{\frac{\alpha x_m}{(\alpha - 1)(2\alpha - 1)}}{\frac{\alpha x_m}{\alpha - 1}} = \frac{\alpha x_m}{(\alpha - 1)(2\alpha - 1)} \cdot \frac{\alpha - 1}{\alpha x_m}
\]

\[
G = \frac{1}{2\alpha - 1}
\]

\begin{tcolorbox}[colback=green!5!white,colframe=resultcolor,title=Result for Part 1]
\[
\boxed{G = \frac{1}{2\alpha - 1}}
\]

This elegant result shows that the Gini coefficient for a Pareto distribution depends only on the shape parameter $\alpha$. As $\alpha$ increases, the inequality (measured by $G$) decreases. For example:
\begin{itemize}
    \item When $\alpha = 1.5$: $G = \frac{1}{2} = 0.5$ (high inequality)
    \item When $\alpha = 2$: $G = \frac{1}{3} \approx 0.333$ (moderate inequality)
    \item When $\alpha = 5$: $G = \frac{1}{9} \approx 0.111$ (low inequality)
\end{itemize}
\end{tcolorbox}

\section{Solution to Part 2: Empirical Gini Coefficient}

\subsection{Step 1: Define the Estimator and Dataset}

For a finite sample $\{x_1, \ldots, x_n\}$, the Gini coefficient is estimated by:
\[
\hat{G} = \frac{\sum_{i=1}^n \sum_{j=1}^n |x_i - x_j|}{2n^2 \bar{x}}
\]

Given dataset: $\{1, 3, 6, 10\}$, so $n = 4$.

\subsection{Step 2: Calculate the Sample Mean}

\[
\bar{x} = \frac{1}{n}\sum_{i=1}^n x_i = \frac{1 + 3 + 6 + 10}{4} = \frac{20}{4} = 5
\]

\subsection{Step 3: Compute All Pairwise Absolute Differences}

We need to calculate $|x_i - x_j|$ for all pairs $(i,j)$. Let's create a difference matrix:

\begin{table}[h]
\centering
\begin{tabular}{c|cccc|c}
\toprule
$|x_i - x_j|$ & $x_1=1$ & $x_2=3$ & $x_3=6$ & $x_4=10$ & Row Sum \\
\midrule
$x_1=1$ & 0 & 2 & 5 & 9 & 16 \\
$x_2=3$ & 2 & 0 & 3 & 7 & 12 \\
$x_3=6$ & 5 & 3 & 0 & 4 & 12 \\
$x_4=10$ & 9 & 7 & 4 & 0 & 20 \\
\midrule
\textbf{Total} & & & & & \textbf{60} \\
\bottomrule
\end{tabular}
\caption{Pairwise absolute differences matrix}
\end{table}

\textbf{Detailed calculations for each row:}

\textbf{Row 1 ($x_1 = 1$):}
\begin{align*}
|1-1| &= 0 \\
|1-3| &= |{-2}| = 2 \\
|1-6| &= |{-5}| = 5 \\
|1-10| &= |{-9}| = 9 \\
\text{Sum} &= 0 + 2 + 5 + 9 = 16
\end{align*}

\textbf{Row 2 ($x_2 = 3$):}
\begin{align*}
|3-1| &= 2 \\
|3-3| &= 0 \\
|3-6| &= |{-3}| = 3 \\
|3-10| &= |{-7}| = 7 \\
\text{Sum} &= 2 + 0 + 3 + 7 = 12
\end{align*}

\textbf{Row 3 ($x_3 = 6$):}
\begin{align*}
|6-1| &= 5 \\
|6-3| &= 3 \\
|6-6| &= 0 \\
|6-10| &= |{-4}| = 4 \\
\text{Sum} &= 5 + 3 + 0 + 4 = 12
\end{align*}

\textbf{Row 4 ($x_4 = 10$):}
\begin{align*}
|10-1| &= 9 \\
|10-3| &= 7 \\
|10-6| &= 4 \\
|10-10| &= 0 \\
\text{Sum} &= 9 + 7 + 4 + 0 = 20
\end{align*}

\textbf{Grand total:}
\[
S = \sum_{i=1}^4 \sum_{j=1}^4 |x_i - x_j| = 16 + 12 + 12 + 20 = 60
\]

\subsection{Step 4: Calculate the Gini Coefficient}

Now we can compute $\hat{G}$:
\[
\hat{G} = \frac{S}{2n^2 \bar{x}} = \frac{60}{2 \times 4^2 \times 5} = \frac{60}{2 \times 16 \times 5} = \frac{60}{160} = \frac{3}{8} = 0.375
\]

\begin{tcolorbox}[colback=green!5!white,colframe=resultcolor,title=Result for Part 2]
\[
\boxed{\hat{G} = 0.375 = \frac{3}{8}}
\]

\textbf{Interpretation:} A Gini coefficient of 0.375 indicates moderate inequality in the dataset. This value is between 0 (perfect equality) and 1 (perfect inequality). The dataset $\{1, 3, 6, 10\}$ shows some dispersion, with the largest value being 10 times the smallest value.
\end{tcolorbox}

\section{Solution to Part 3: Bounds on the Gini Coefficient}

\subsection{Proof that $0 \leq \hat{G} \leq 1$ for Non-Negative Data}

\textbf{Theorem:} For any non-negative dataset $\{x_1, x_2, \ldots, x_n\}$ with $x_i \geq 0$ for all $i$, the empirical Gini coefficient satisfies:
\[
0 \leq \hat{G} \leq 1
\]

\subsection{Part A: Lower Bound ($\hat{G} \geq 0$)}

The Gini coefficient is defined as:
\[
\hat{G} = \frac{\sum_{i=1}^n \sum_{j=1}^n |x_i - x_j|}{2n^2 \bar{x}}
\]

\textbf{Step 1:} Since absolute values are always non-negative:
\[
|x_i - x_j| \geq 0 \quad \text{for all } i, j
\]

Therefore:
\[
\sum_{i=1}^n \sum_{j=1}^n |x_i - x_j| \geq 0
\]

\textbf{Step 2:} The denominator is positive when at least one $x_i > 0$ (otherwise $\bar{x} = 0$ and $\hat{G}$ is undefined):
\[
2n^2 \bar{x} = 2n^2 \cdot \frac{\sum_{i=1}^n x_i}{n} = 2n \sum_{i=1}^n x_i > 0
\]

\textbf{Step 3:} Since both numerator and denominator are non-negative:
\[
\hat{G} = \frac{\text{(non-negative)}}{\text{(positive)}} \geq 0
\]

\textbf{When does $\hat{G} = 0$?} Equality holds when:
\[
\sum_{i=1}^n \sum_{j=1}^n |x_i - x_j| = 0
\]
This occurs if and only if $x_i = x_j$ for all $i, j$, meaning all values are equal (perfect equality).

\subsection{Part B: Upper Bound ($\hat{G} \leq 1$)}

\textbf{Key inequality:} For non-negative values $x_i, x_j \geq 0$, we have:
\[
|x_i - x_j| \leq x_i + x_j
\]

\textbf{Proof of key inequality:} Without loss of generality, assume $x_i \geq x_j \geq 0$. Then:
\begin{align*}
|x_i - x_j| &= x_i - x_j \\
&\leq x_i - x_j + 2x_j \quad \text{(since } x_j \geq 0\text{)} \\
&= x_i + x_j
\end{align*}

\textbf{Step 1:} Applying this inequality to the double sum:
\[
\sum_{i=1}^n \sum_{j=1}^n |x_i - x_j| \leq \sum_{i=1}^n \sum_{j=1}^n (x_i + x_j)
\]

\textbf{Step 2:} Expand the right-hand side:
\[
\sum_{i=1}^n \sum_{j=1}^n (x_i + x_j) = \sum_{i=1}^n \sum_{j=1}^n x_i + \sum_{i=1}^n \sum_{j=1}^n x_j
\]

For the first term, $x_i$ is constant with respect to the inner sum over $j$:
\[
\sum_{i=1}^n \sum_{j=1}^n x_i = \sum_{i=1}^n x_i \sum_{j=1}^n 1 = \sum_{i=1}^n x_i \cdot n = n \sum_{i=1}^n x_i
\]

Similarly for the second term:
\[
\sum_{i=1}^n \sum_{j=1}^n x_j = \sum_{i=1}^n \left(\sum_{j=1}^n x_j\right) = n \sum_{j=1}^n x_j
\]

Therefore:
\[
\sum_{i=1}^n \sum_{j=1}^n (x_i + x_j) = n \sum_{i=1}^n x_i + n \sum_{j=1}^n x_j = 2n \sum_{i=1}^n x_i
\]

\textbf{Step 3:} Substitute into the Gini formula:
\[
\hat{G} = \frac{\sum_{i=1}^n \sum_{j=1}^n |x_i - x_j|}{2n^2 \bar{x}} \leq \frac{2n \sum_{i=1}^n x_i}{2n^2 \bar{x}}
\]

\textbf{Step 4:} Recall that $\bar{x} = \frac{1}{n}\sum_{i=1}^n x_i$, so:
\[
\hat{G} \leq \frac{2n \sum_{i=1}^n x_i}{2n^2 \cdot \frac{1}{n}\sum_{i=1}^n x_i} = \frac{2n \sum_{i=1}^n x_i}{2n \sum_{i=1}^n x_i} = 1
\]

\subsection{When does $\hat{G} = 1$?}

Maximum inequality occurs when $|x_i - x_j| = x_i + x_j$ for all pairs. From our inequality proof, this happens when one of the values equals zero.

\textbf{Example:} Consider the dataset $\{0, 0, \ldots, 0, a\}$ where $n-1$ values are zero and one value is $a > 0$.

Then $\bar{x} = \frac{a}{n}$, and:
\begin{align*}
\sum_{i=1}^n \sum_{j=1}^n |x_i - x_j| &= 2(n-1) \times a = 2(n-1)a
\end{align*}

(Each of the $n-1$ zeros paired with $a$ contributes $|0-a| = a$, counted twice due to symmetry)

Therefore:
\[
\hat{G} = \frac{2(n-1)a}{2n^2 \cdot \frac{a}{n}} = \frac{2(n-1)a}{2na} = \frac{n-1}{n} = 1 - \frac{1}{n} \to 1 \text{ as } n \to \infty
\]

For finite $n$, the maximum is $\frac{n-1}{n}$, which approaches 1 as the sample size increases.

\begin{tcolorbox}[colback=green!5!white,colframe=resultcolor,title=Result for Part 3]
\textbf{Proven:}
\[
\boxed{0 \leq \hat{G} \leq 1}
\]

\textbf{Summary:}
\begin{itemize}
    \item $\hat{G} = 0$ occurs when all values are equal (perfect equality)
    \item $\hat{G} = 1$ (or close to 1) occurs when all wealth is concentrated in one observation (perfect inequality)
    \item The bounds hold for any non-negative dataset
\end{itemize}
\end{tcolorbox}

\section{Summary and Key Insights}

\begin{tcolorbox}[colback=yellow!10!white,colframe=orange,title=Final Summary]
\subsection*{Main Results}

\textbf{1. Theoretical Gini for Pareto Distribution:}
\[
G = \frac{1}{2\alpha - 1}
\]
This shows that higher shape parameter $\alpha$ leads to lower inequality.

\textbf{2. Empirical Gini for Dataset $\{1, 3, 6, 10\}$:}
\[
\hat{G} = 0.375
\]
This indicates moderate inequality in the sample.

\textbf{3. Universal Bounds:}
\[
0 \leq \hat{G} \leq 1 \quad \text{for all non-negative datasets}
\]

\subsection*{Applications}

The Gini coefficient is widely used in:
\begin{itemize}
    \item \textbf{Economics:} Measuring income and wealth inequality
    \item \textbf{Healthcare:} Assessing disparities in health outcomes
    \item \textbf{Ecology:} Measuring biodiversity and species distribution
    \item \textbf{Machine Learning:} Feature selection and decision tree splitting criteria
\end{itemize}

\subsection*{Key Properties}

\begin{enumerate}
    \item \textbf{Scale invariance:} Multiplying all values by a constant doesn't change $G$
    \item \textbf{Translation sensitivity:} Adding a constant to all values changes $G$
    \item \textbf{Population independence:} For Pareto, $G$ depends only on $\alpha$, not $x_m$
    \item \textbf{Interpretability:} Values between 0.25-0.35 indicate reasonable equality, above 0.50 indicates high inequality
\end{enumerate}
\end{tcolorbox}

\section*{Conclusion}

This problem demonstrates the mathematical rigor required to analyze inequality measures. The Pareto distribution's analytical Gini coefficient provides a theoretical foundation, while the empirical estimate shows practical application. The proof of bounds ensures the measure is well-behaved and interpretable across all datasets.

The step-by-step derivations reveal the deep connections between probability theory, calculus, and statistical measurement, showcasing how abstract mathematical concepts translate into practical tools for understanding real-world phenomena.

\end{document}
